\chapter{Gravity from Entropy Gradients}
\label{chap:gravity}

\section{Motivation: Why Emergence?}

Classical general relativity models the dynamics of spacetime geometry by
relating curvature to stress--energy via the Einstein field equations.
Lamphrodynamics modifies this paradigm by introducing an entropy--driven
correction term arising from the lamphrodynamic fields
$\Psi := (\PhiField,\vField,\SField)$ defined in Volume~I.

Informally, the guiding principle of \rsvpabbr{} is that gravitational
attraction is not a fundamental interaction mediated by a metric tensor, but is
an \emph{entropic phenomenon} induced by the dynamics of $\Psi$. In this view
spacetime geometry emerges as an effective description of deeper lamphrodynamic
degrees of freedom. The aim of this chapter is to explain, derive, and justify
this perspective in a mathematically precise manner.

From a physical standpoint, this interpretation is motivated by the observation
that regions of high entropy density tend to induce lamphrodynamic flows which
pull field configurations into lower--energy, higher--coherence states. The
associated effective stress--energy contributes an additional curvature
term in the field equations, one which reduces to the usual Einstein tensor in
the limit where entropy gradients vanish.

\section{Review of Lamphrodynamic Variational Structure}

Recall from Volume~I that lamphrodynamic field configurations are described by
the action functional
\[
  \mathcal{S}_{\mathrm{AKSZ}}[\Psi]
  = \int_M \omega_{-1}(\Psi, \mathrm{d}\Psi) + \mathcal{V}[\Psi],
\]
constructed via the AKSZ procedure with target the $(-1)$–shifted symplectic
stack $\mathcal{P}$ introduced in \cref{chap:shifted-symplectic}. The Euler--%
Lagrange equations take the schematic form
\[
  E_\Psi(\PhiField,\vField,\SField) = 0,
\]
coupled to the metric variation arising from the underlying AKSZ action. We
now recall the basic structure needed for the gravitational sector.

\begin{theorem}[Metric Variation of the AKSZ--RSVP Action]
\label{thm:metric-variation}
Let $(M,g)$ be a lamphrodynamic spacetime and $\Psi=(\PhiField,\vField,\SField)$
a smooth configuration. Then the variational derivative of the AKSZ--RSVP action
with respect to $g$ is given by
\[
  \frac{\delta \mathcal{S}_{\mathrm{AKSZ}}}{\delta g^{\mu\nu}}
  = -\frac{1}{2}\sqrt{|g|}\,
    \bigl( T_{\mu\nu}^{\mathrm{matter}} + \Theta_{\mu\nu} \bigr),
\]
where $\Theta_{\mu\nu}$ is an entropy--induced tensor determined by
$(\PhiField,\vField,\SField)$ and the $(-1)$–shifted symplectic form
$\omega_{-1}$.
\end{theorem}

\begin{proof}[Sketch of proof]
The variation follows from the AKSZ construction and the explicit form of the
target $(-1)$–shifted symplectic structure. The term $\Theta_{\mu\nu}$ arises
from the nontrivial dependence of $\omega_{-1}$ and $\mathcal{V}$ on the metric
through entropy gradients and lamphrodynamic flow. The full derivation is given
in \cref{sec:explicit-theta}.
\end{proof}

\section{Derivation of Modified Einstein Equations}

Combining \cref{thm:metric-variation} with the standard variation of the
Einstein--Hilbert action leads to the modified Einstein field equations
\begin{equation}
  G_{\mu\nu} + \Theta_{\mu\nu}
  = 8\pi G\, T_{\mu\nu}^{\mathrm{matter}},
  \label{eq:modified-einstein}
\end{equation}
where $\Theta_{\mu\nu}$ depends on $(\PhiField,\vField,\SField)$ and vanishes
when the lamphrodynamic fields are trivial. By construction,
$\Theta_{\mu\nu}$ is symmetric and covariantly conserved:
\[
  \nabla^\mu \Theta_{\mu\nu} = 0,
\]
which follows directly from the $(-1)$–shifted symplectic structure and the
AKSZ master equation.

Equation~\eqref{eq:modified-einstein} should be interpreted as a set of
effective gravitational field equations, where curvature is driven not only by
ordinary matter but also by entropy dynamics. The geometry is therefore
emergent in the sense that it is determined by lamphrodynamic evolution rather
than postulated a priori.

\section{Explicit Computation of the Entropic Correction Tensor}

We now give an explicit formula for the tensor $\Theta_{\mu\nu}$ introduced in
\cref{thm:metric-variation}. As in Volume~I, let
$\Psi = (\PhiField,\vField,\SField)$ denote the lamphrodynamic fields, and
write $J^\mu_S$ for the entropy--flux vector. The $(-1)$–shifted symplectic
structure supplies a natural bilinear pairing
\[
  \omega_{-1} : T_{-1}\mathcal{P} \times T_{-1}\mathcal{P} \longrightarrow \mathbb{R}[-1],
\]
whose local expression, in the gauge chosen in \cref{chap:aksz-rsvp}, takes the
schematic form
\[
  \omega_{-1}(\Psi,\mathrm{d}\Psi)
   = \SField\, g_{\mu\nu}\,\mathrm{d}\PhiField \wedge \mathrm{d}v^\nu
   + \cdots .
\]
Varying this with respect to the metric and extracting the $\delta g^{\mu\nu}$
component produces the following.

\begin{proposition}
\label{prop:theta-explicit}
The entropic correction tensor associated to the lamphrodynamic fields is given
in local coordinates by
\begin{equation}
\label{eq:theta-explicit}
  \Theta_{\mu\nu}
  := \alpha\,\Bigl(
        \nabla_\mu \SField\,\nabla_\nu\PhiField
      + \nabla_\mu \PhiField\,\nabla_\nu\SField
      - g_{\mu\nu}\, \nabla_\rho\SField\,\nabla^\rho\PhiField
     \Bigr)
  + \beta\,\Bigl(
        J^S_\mu\,J^S_\nu
      - \tfrac{1}{2} g_{\mu\nu} J^S_\rho J_S^\rho
     \Bigr),
\end{equation}
where $\alpha,\beta$ are universal coupling constants determined by the AKSZ
target and the normalization of $\omega_{-1}$.
\end{proposition}

\begin{proof}
The term proportional to $\alpha$ arises from metric variations of the
$(-1)$–shifted symplectic pairing between $\PhiField$ and $\SField$; the term
proportional to $\beta$ arises from the kinetic part of the entropy--flux
sector. The full derivation follows the standard variation procedure for AKSZ
actions, using the explicit form of $\omega_{-1}$ in a local coordinate chart
and the master equation to guarantee symmetry and conservation. Details appear
in Appendix~\ref{app:aksz-variation-details}.
\end{proof}

\begin{remark}
The tensor $\Theta_{\mu\nu}$ plays the role of an ``entropic stress--energy
tensor,'' but unlike the usual perfect–fluid form, it captures both entropy
gradients (through $\alpha$) and entropy fluxes (through $\beta$).
\end{remark}

\subsection{Weak--Field Limit and Newtonian Approximation}

To understand the physical content of \eqref{eq:modified-einstein} and
\eqref{eq:theta-explicit}, we linearize the metric around Minkowski space.
Write $g_{\mu\nu} = \eta_{\mu\nu} + h_{\mu\nu}$ with $|h_{\mu\nu}|\ll 1$.
Retaining only linear terms in $h_{\mu\nu}$ and the lamphrodynamic fields
yields the approximate field equations
\begin{equation}
\label{eq:weak-field}
  \Box h_{\mu\nu}
   = -16\pi G\,
      \Bigl(
          T^{\mathrm{matter}}_{\mu\nu}
        + \Theta^{(1)}_{\mu\nu}
      \Bigr),
\end{equation}
where $\Theta^{(1)}_{\mu\nu}$ denotes the linear part of
\eqref{eq:theta-explicit}. In the Newtonian gauge and for slowly varying
fields, this reduces to the effective Poisson equation
\begin{equation}
\label{eq:poisson}
  \Delta \Phi_{\mathrm{grav}}
  = 4\pi G\bigl(\rho_{\mathrm{matter}}
             + \rho_{\mathrm{entropy}}\bigr),
\end{equation}
with $\rho_{\mathrm{entropy}}$ obtained from the time--time component of
$\Theta^{(1)}_{\mu\nu}$.

This shows that entropy gradients and entropy fluxes contribute additional
``mass–density'' terms, leading to enhanced gravitational attraction in regions
with strong lamphrodynamic activity.

\subsection{Comparison with General Relativity}

In the limit $\alpha,\beta\to 0$ (or equivalently when $\SField$ and $J_S$
vanish), equation~\eqref{eq:poisson} reduces to the classical Newtonian
potential equation with source $\rho_{\mathrm{matter}}$ only. Hence general
relativity is recovered as a degenerate limit of the lamphrodynamic theory.

By contrast, when $\SField$ and $J_S$ are non--trivial, the right--hand side of
\eqref{eq:poisson} includes additional contributions, providing a potential
explanation for phenomena traditionally attributed to dark matter or modified
gravity. Subsequent chapters will analyze these possibilities in detail.

\subsection{Phenomenological Interpretation of $\rho_{\mathrm{entropy}}$}

Equation~\eqref{eq:poisson} suggests interpreting
\[
  \rho_{\mathrm{entropy}}
   := \frac{1}{4\pi G}\,\Theta^{(1)}_{00}
\]
as an effective mass--density associated to lamphrodynamic degrees of freedom.
Two qualitatively distinct mechanisms appear:

\begin{enumerate}
\item \emph{Static entropic gradients}: regions with $\nabla_i\SField\neq0$
produce a contribution to $\rho_{\mathrm{entropy}}$ similar in form to scalar
field sources, but with reversed sign depending on the parameters
$(\alpha,\beta)$. In typical lamphrodynamic regimes, $\alpha > 0$ yields
additional attractive potential.

\item \emph{Dynamical entropy flux}: regions with $J^S_i\neq0$
contribute via the flux--flux term in \eqref{eq:theta-explicit}, producing
attractive gravitational effects even when $\SField$ is spatially
homogeneous.
\end{enumerate}

These two effects are conceptually different: entropy density gradients
correspond to spatial variation in the thermodynamic state of the plenum,
whereas entropy flux corresponds to ongoing irreversible processes or
negentropic flows. Their coexistence in $\Theta_{\mu\nu}$ is a distinctive
signature of the lamphrodynamic theory absent in general relativity.

\subsection{Observational Consequences}

In the weak--field, quasi--stationary regime, the effective density
$\rho_{\mathrm{entropy}}$ acts as an additional gravitational source.
This leads to several immediate observational consequences:

\begin{enumerate}
\item \emph{Galaxy rotation curves.} A spatially varying $\SField$ can modify
the radial gravitational potential without invoking dark matter halos.
Preliminary rotation--curve fits are discussed in \cref{chap:dark-matter}.

\item \emph{Cluster binding.} Enhanced effective mass at large scales may
stabilize galaxy clusters consistent with observed velocity dispersions.

\item \emph{Gravitational lensing.} The entropy--flux term modifies the
deflection angle for light rays, potentially explaining lensing anomalies in
clusters or strong--lensing systems.
\end{enumerate}

These consequences will be quantified and compared with observational data in
\cref{part:observational}.

\subsection{Dimensional Analysis and Coupling Constants}

The constants $\alpha,\beta$ appearing in \eqref{eq:theta-explicit} are
dimensionful and depend on the normalization of the AKSZ target, as well as
the units in which entropy and potential are measured. A convenient choice is
\[
  [\PhiField] = \text{energy per unit mass},\qquad
  [\SField] = k_B\,(\text{dimensionless}),
\]
so that $(\alpha,\beta)$ carry dimensions of length$^2$ in geometrized units.
In the classical limit $\hbar\to0$ (where the AKSZ structure becomes purely
classical), one may treat $\alpha,\beta$ as phenomenological constants to be
fit by observation; in the quantum regime, they are determined by the target
geometry of the BV theory.

\subsection{Recovery of Newtonian Gravity}

Setting $\SField=0$ and $J_S=0$ exactly, the tensor $\Theta_{\mu\nu}$ vanishes
identically, and \eqref{eq:poisson} reduces to
\[
  \Delta \Phi_{\mathrm{grav}} = 4\pi G \rho_{\mathrm{matter}},
\]
the usual Newtonian gravity. Thus, the lamphrodynamic theory strictly contains
general relativity as a special case, without introducing additional fields
beyond $(\PhiField,\vField,\SField)$.

\subsection{Relation to Scalar--Tensor Theories}

At first sight, $\Theta_{\mu\nu}$ resembles the stress--energy tensor of a
scalar field $\varphi$ in scalar--tensor gravity. However, the presence of
$\vField$ and the flux sector $J_S$ distinguishes the lamphrodynamic case:
there is no single scalar field but a coupled system of scalar potential,
vector flow, and entropy density. Moreover, $\SField$ is thermodynamic rather
than fundamental, encoding local entropy rather than a new fundamental
scalar. This qualitative difference will be crucial in subsequent chapters,
especially when comparing with scalar--tensor screening mechanisms.

\subsection{Summary of the Weak--Field Analysis}

The main conclusions of this chapter are:
\begin{itemize}
\item the lamphrodynamic correction tensor $\Theta_{\mu\nu}$ provides an
effective gravitational source proportional to entropy gradients and fluxes;
\item these corrections produce additional gravitational attraction in
situations with nontrivial lamphrodynamic activity;
\item Newtonian gravity is recovered when $\SField$ and $J_S$ vanish or are
negligible;
\item the resulting phenomenology can potentially address astrophysical
anomalies without invoking nonbaryonic dark matter.
\end{itemize}

In the next chapter we compare these predictions with other emergent--gravity
frameworks, especially Jacobson's thermodynamic derivation of Einstein's
equations and Verlinde's entropic gravity proposals, highlighting both
conceptual and quantitative differences.


